% Input fragment; use the Chinese Volume 0 preamble.
% Assign a unique label before integration; remove unused optional blocks.
\chapter{待定章节标题}\label{template:v0}
\begin{chaptersummary}
【作者待填写：约 100--250 字，连接已有结果、剩余问题、本章选择、
后续用途与一个核心思想。内容以正文为依据。】
\end{chaptersummary}
\learninggoals{【填写 2--4 个可检验目标，如构造、推导、比较或验证。】}
\paragraph{核心问题}【此前结果尚未解决的具体问题是什么？】

\section{问题来源与适用条件}
【交代离散算子的来源、矩阵类别、变量、尺度及所需假设。】

\section{关键关系与最小例子}
【在此加入已确认的数学骨架；公式前说明用途，公式后解释结果。】
【填写一个可手算或可复现的最小例子及验收指标。】

% Optional: retain only for a genuinely missing conceptual motivation.
\conceptbridge{【填写读者可能提出的具体疑问。】}
{【作者待补充：80--300 字，解释为什么需要这一步及选择依据。】}
{【已有对象】$\longrightarrow$【新对象】}

\section{成本与计算接口}
\costnote{【填写运算、存储或通信成本的计数口径、适用条件及比较指标。】}
% Optional: \coderef{existing/relative_code_path.py}
% Optional: \crossref{A specific result needed from another chapter or volume.}
\aiexercise{【指定符号、假设、推导起点与可检查的计算或例子。】}
{【改变一个参数或假设，提出能检验机制的后续问题。】}
% Volume 0 studyroute takes TWO arguments.
\studyroute{【明确略去哪些证明、代数或实现细节，并添加已核实的文献。】}
{【填写精确、通用的英文检索词】}
